% =====================================================================
% style/preamble.tex - Preambolo condiviso del libro di armonia
% =====================================================================
% Tracciato (pubblico): contiene SOLO la configurazione tipografica e i
% pacchetti, nessun contenuto del libro. Caricato dai file principali con
% % =====================================================================
% style/preamble.tex - Preambolo condiviso del libro di armonia
% =====================================================================
% Tracciato (pubblico): contiene SOLO la configurazione tipografica e i
% pacchetti, nessun contenuto del libro. Caricato dai file principali con
% % =====================================================================
% style/preamble.tex - Preambolo condiviso del libro di armonia
% =====================================================================
% Tracciato (pubblico): contiene SOLO la configurazione tipografica e i
% pacchetti, nessun contenuto del libro. Caricato dai file principali con
% % =====================================================================
% style/preamble.tex - Preambolo condiviso del libro di armonia
% =====================================================================
% Tracciato (pubblico): contiene SOLO la configurazione tipografica e i
% pacchetti, nessun contenuto del libro. Caricato dai file principali con
% \input{preamble} subito dopo \documentclass{memoir} (LuaLaTeX).
% L'engine e' fissato in .latexmkrc (lualatex). Vedi .claude/context/STACK.md.
% ---------------------------------------------------------------------

% --- Font e Unicode (LuaLaTeX) ---
\usepackage{fontspec}
\usepackage{unicode-math}
\setmainfont{Libertinus Serif}
\setsansfont{Libertinus Sans}
\setmathfont{Libertinus Math}

% --- Lingua italiana ---
\usepackage[italian]{babel}

% --- Microtipografia editoriale (espansione + protrusione) ---
\usepackage[final,expansion=true,protrusion=true]{microtype}

% --- Virgolette coerenti con la lingua ---
\usepackage[autostyle=true,italian=guillemets]{csquotes}

% --- Matematica (intervalli, rapporti, formule armoniche) ---
\usepackage{amsmath}
\usepackage{mathtools}

% --- Tabelle ---
\usepackage{booktabs}
\usepackage{tabularx}

% --- Figure e disegni vettoriali (cerchio delle quinte, schemi tonali) ---
\usepackage{graphicx}
\usepackage{float}
\usepackage[font=small,labelfont=bf]{caption}
\usepackage{xcolor}
\usepackage{tikz}
\usepackage{pgfplots}
\pgfplotsset{compat=1.18}

% --- Bibliografia (biblatex + biber); ogni main aggiunge \addbibresource ---
\usepackage[backend=biber,style=authoryear,language=italian]{biblatex}

% --- Indice analitico e glossario dei termini ---
\usepackage{imakeidx}
\makeindex[intoc]
\usepackage[acronym,toc]{glossaries}
\makeglossaries

% --- Link e riferimenti (hyperref penultimo, cleveref ultimo) ---
\usepackage[hidelinks,unicode]{hyperref}
\usepackage{cleveref}

% --- Macro proprie di notazione armonica ---
\usepackage{harmony-macros}
 subito dopo \documentclass{memoir} (LuaLaTeX).
% L'engine e' fissato in .latexmkrc (lualatex). Vedi .claude/context/STACK.md.
% ---------------------------------------------------------------------

% --- Font e Unicode (LuaLaTeX) ---
\usepackage{fontspec}
\usepackage{unicode-math}
\setmainfont{Libertinus Serif}
\setsansfont{Libertinus Sans}
\setmathfont{Libertinus Math}

% --- Lingua italiana ---
\usepackage[italian]{babel}

% --- Microtipografia editoriale (espansione + protrusione) ---
\usepackage[final,expansion=true,protrusion=true]{microtype}

% --- Virgolette coerenti con la lingua ---
\usepackage[autostyle=true,italian=guillemets]{csquotes}

% --- Matematica (intervalli, rapporti, formule armoniche) ---
\usepackage{amsmath}
\usepackage{mathtools}

% --- Tabelle ---
\usepackage{booktabs}
\usepackage{tabularx}

% --- Figure e disegni vettoriali (cerchio delle quinte, schemi tonali) ---
\usepackage{graphicx}
\usepackage{float}
\usepackage[font=small,labelfont=bf]{caption}
\usepackage{xcolor}
\usepackage{tikz}
\usepackage{pgfplots}
\pgfplotsset{compat=1.18}

% --- Bibliografia (biblatex + biber); ogni main aggiunge \addbibresource ---
\usepackage[backend=biber,style=authoryear,language=italian]{biblatex}

% --- Indice analitico e glossario dei termini ---
\usepackage{imakeidx}
\makeindex[intoc]
\usepackage[acronym,toc]{glossaries}
\makeglossaries

% --- Link e riferimenti (hyperref penultimo, cleveref ultimo) ---
\usepackage[hidelinks,unicode]{hyperref}
\usepackage{cleveref}

% --- Macro proprie di notazione armonica ---
\usepackage{harmony-macros}
 subito dopo \documentclass{memoir} (LuaLaTeX).
% L'engine e' fissato in .latexmkrc (lualatex). Vedi .claude/context/STACK.md.
% ---------------------------------------------------------------------

% --- Font e Unicode (LuaLaTeX) ---
\usepackage{fontspec}
\usepackage{unicode-math}
\setmainfont{Libertinus Serif}
\setsansfont{Libertinus Sans}
\setmathfont{Libertinus Math}

% --- Lingua italiana ---
\usepackage[italian]{babel}

% --- Microtipografia editoriale (espansione + protrusione) ---
\usepackage[final,expansion=true,protrusion=true]{microtype}

% --- Virgolette coerenti con la lingua ---
\usepackage[autostyle=true,italian=guillemets]{csquotes}

% --- Matematica (intervalli, rapporti, formule armoniche) ---
\usepackage{amsmath}
\usepackage{mathtools}

% --- Tabelle ---
\usepackage{booktabs}
\usepackage{tabularx}

% --- Figure e disegni vettoriali (cerchio delle quinte, schemi tonali) ---
\usepackage{graphicx}
\usepackage{float}
\usepackage[font=small,labelfont=bf]{caption}
\usepackage{xcolor}
\usepackage{tikz}
\usepackage{pgfplots}
\pgfplotsset{compat=1.18}

% --- Bibliografia (biblatex + biber); ogni main aggiunge \addbibresource ---
\usepackage[backend=biber,style=authoryear,language=italian]{biblatex}

% --- Indice analitico e glossario dei termini ---
\usepackage{imakeidx}
\makeindex[intoc]
\usepackage[acronym,toc]{glossaries}
\makeglossaries

% --- Link e riferimenti (hyperref penultimo, cleveref ultimo) ---
\usepackage[hidelinks,unicode]{hyperref}
\usepackage{cleveref}

% --- Macro proprie di notazione armonica ---
\usepackage{harmony-macros}
 subito dopo \documentclass{memoir} (LuaLaTeX).
% L'engine e' fissato in .latexmkrc (lualatex). Vedi .claude/context/STACK.md.
% ---------------------------------------------------------------------

% --- Font e Unicode (LuaLaTeX) ---
\usepackage{fontspec}
\usepackage{unicode-math}
\setmainfont{Libertinus Serif}
\setsansfont{Libertinus Sans}
\setmathfont{Libertinus Math}

% --- Lingua italiana ---
\usepackage[italian]{babel}

% --- Microtipografia editoriale (espansione + protrusione) ---
\usepackage[final,expansion=true,protrusion=true]{microtype}

% --- Virgolette coerenti con la lingua ---
\usepackage[autostyle=true,italian=guillemets]{csquotes}

% --- Matematica (intervalli, rapporti, formule armoniche) ---
\usepackage{amsmath}
\usepackage{mathtools}

% --- Tabelle ---
\usepackage{booktabs}
\usepackage{tabularx}

% --- Figure e disegni vettoriali (cerchio delle quinte, schemi tonali) ---
\usepackage{graphicx}
\usepackage{float}
\usepackage[font=small,labelfont=bf]{caption}
\usepackage{xcolor}
\usepackage{tikz}
\usepackage{pgfplots}
\pgfplotsset{compat=1.18}

% --- Bibliografia (biblatex + biber); ogni main aggiunge \addbibresource ---
\usepackage[backend=biber,style=authoryear,language=italian]{biblatex}

% --- Indice analitico e glossario dei termini ---
\usepackage{imakeidx}
\makeindex[intoc]
\usepackage[acronym,toc]{glossaries}
\makeglossaries

% --- Link e riferimenti (hyperref penultimo, cleveref ultimo) ---
\usepackage[hidelinks,unicode]{hyperref}
\usepackage{cleveref}

% --- Macro proprie di notazione armonica ---
\usepackage{harmony-macros}
